\BookChapter{8}

\section[第8章\q 豪车悖论]{第8章\\ 豪车悖论}

{\kaishu

其实别人不会像你那样在意你有多少财产。

}

泊车员这份工作最棒的一点就是可以驾驶顶级豪车，毕竟客人们都是开着法拉利、兰博基尼、劳斯莱斯来的——如同贵族的舰队。

我曾经的梦想就是拥有一辆这样的车，因为（那时我认为）开豪车会让别人觉得你是个成功人士：聪明，富有，有品位，是一个举足轻重的人物。看看我多厉害啊。

讽刺的是，我几乎从来没有注意到是谁坐在这些豪车的驾驶座上。

当你看见有人开着一辆漂亮的车子时，你很少会想“哇，开那辆车的家伙真酷”，相反，你只会想“哇，如果那辆车是我的，别人一定会觉得我很酷”。无论这种念头是有意识的还是只存在于潜意识中，人们在这种情况下的态度就是如此。

所以，这里存在一个悖论：我们都想通过财富来告诉其他人，自己应该受到他们的爱慕与敬仰。但事实上，其他人常常会跳过敬仰你这一步。这并不是因为他们觉得你的财富不值得羡慕，而是因为他们会把你的财富当作标尺，转而表达自己渴望被爱慕与敬仰的愿望。

在我儿子出生以后，我给他写了一封信。信中写道：“你或许觉得你需要一辆昂贵的车子、一块豪华的手表和一座很大的房子，但我想告诉你的是，你并非真想得到这些东西本身。你真想得到的是来自他人的尊重和羡慕。你觉得拥有昂贵的东西会让别人尊重和羡慕你，但可惜，别人不会——尤其是那些你希望得到其尊重和羡慕的人。”

在做泊车员时，我在对那些开着法拉利到酒店、享受我艳羡目光的人们的思考中明白了这个道理。在他们所到之处，人们都会这样盯着他们。我相信他们一定很喜欢这种感觉。我敢肯定他们觉得自己获得了别人的仰慕。

他们可曾知道，我其实并不关心他们本人，甚至根本没注意到他们？他们可曾知道，我注意到的只是他们的车，心里想的只有我自己坐在驾驶座上的样子？

他们在买法拉利的时候是否盘算着获得别人的仰慕，而没有意识到我——以及其他很多人——会为他们的豪车着迷，却一秒也不会考虑到坐在驾驶座上的人——他们本身？

同样的道理是否也适用于那些住在豪宅里的人？答案几乎是肯定的。

那些佩戴昂贵珠宝、穿着华丽服饰的人呢？答案同样如此。

在此，我的重点不是让你放弃追求财富或豪车的梦想。这二者依然是我喜欢的。

但我们需要知道，虽然人们的确都渴望得到他人的尊重和羡慕，但通过用金钱购买昂贵之物获得的尊重和羡慕可能远比你想象中少。如果获得尊重和羡慕是你的目标，那么一定要注意选择正确的方法。比起豪车，谦虚、善良和同情心等人格特质才能帮你获得更多尊重。

关于法拉利的问题我还没讲完。下一章我会讲讲与“豪车悖论”相关的另一个问题。
